\documentclass[a4paper,10pt]{article}
\usepackage[utf8]{inputenc}
\usepackage[T1]{fontenc}
\usepackage[french]{babel}
\usepackage[margin=1.5cm]{geometry}
\usepackage{xcolor}
\usepackage{tcolorbox}
\usepackage{enumitem}
\usepackage{multicol}
\usepackage{lmodern}
\usepackage{tabularx}
\usepackage{booktabs}
\usepackage{listings}

\renewcommand{\familydefault}{\sfdefault}
\pagestyle{empty}
\setlist{nosep, leftmargin=1.5em}

\definecolor{navy}{RGB}{0,51,102}
\definecolor{teal}{RGB}{0,120,120}
\definecolor{crimson}{RGB}{180,0,0}
\definecolor{green}{RGB}{0,120,50}
\definecolor{lightbg}{RGB}{245,248,252}
\definecolor{orange}{RGB}{200,100,0}
\definecolor{codebg}{RGB}{30,30,30}
\definecolor{codetext}{RGB}{220,220,220}

\lstset{
  backgroundcolor=\color{codebg},
  basicstyle=\ttfamily\small\color{codetext},
  breaklines=true, frame=none,
  keywordstyle=\color{cyan},
  commentstyle=\color{gray},
  stringstyle=\color{yellow},
}

\tcbset{basebox/.style={
  colback=lightbg, colframe=navy, coltitle=white,
  fonttitle=\bfseries, arc=2mm, boxrule=0.5mm,
  left=3mm, right=3mm, top=2mm, bottom=2mm, after skip=6pt
}}

\newcommand{\key}[1]{\textbf{\textcolor{navy}{#1}}}
\newcommand{\warn}[1]{\textbf{\textcolor{crimson}{#1}}}
\newcommand{\cmd}[1]{\texttt{\textcolor{navy}{#1}}}

\begin{document}

\begin{center}
  \colorbox{navy}{\parbox{\dimexpr\textwidth-2\fboxsep}{
    \centering\vspace{3pt}
    {\LARGE\bfseries\textcolor{white}{ADMINISTRATION LINUX \& SERVEUR — Fiche de révision}}\\
    {\normalsize\textcolor{white}{BTS SIO — Maël Mignard}}
    \vspace{3pt}
  }}
\end{center}
\vspace{6pt}

\begin{multicols}{2}

% ============================================================
\begin{tcolorbox}[basebox, colframe=navy, title=1. Structure d'une commande Bash]
Syntaxe : \cmd{programme [options] [arguments]} \\[4pt]
\begin{tabular}{ll}
\textbf{Élément} & \textbf{Exemple} \\
\hline
Programme & \cmd{ls}, \cmd{cp}, \cmd{mkdir} \\
Option & \cmd{-l}, \cmd{-a}, \cmd{-h}, \cmd{-rf} \\
Argument & \cmd{/etc/apache2}, \cmd{fichier.txt} \\
\end{tabular}
\vspace{4pt}\\
Ex: \cmd{ls -lah /var/log}\\
\tip{Les options peuvent souvent s'interchanger avec les arguments.}
\end{tcolorbox}

% ============================================================
\begin{tcolorbox}[basebox, colframe=green, title=2. Commandes essentielles]
\begin{tabularx}{\columnwidth}{lX}
\toprule
\textbf{Commande} & \textbf{Description} \\
\midrule
\cmd{ls -lah}    & Liste fichiers (droits, taille lisible, cachés) \\
\cmd{pwd}        & Chemin courant absolu \\
\cmd{cd /path}   & Changer de répertoire \\
\cmd{mkdir -p}   & Créer répertoire (parents inclus) \\
\cmd{touch file} & Créer fichier vide \\
\cmd{cat file}   & Afficher contenu \\
\cmd{nano file}  & Éditer fichier (terminal) \\
\cmd{rm -rf dir} & Supprimer récursivement (\warn{dangereux}) \\
\cmd{cp -R src dst} & Copier récursivement \\
\cmd{mv src dst} & Déplacer / renommer \\
\cmd{scp user@ip:/src /dst} & Copie via SSH \\
\cmd{grep "mot" fichier} & Recherche de texte \\
\cmd{find . -name "*.txt"} & Recherche de fichier \\
\cmd{ps aux}     & Processus en cours \\
\cmd{top / htop} & Moniteur système interactif \\
\cmd{kill -9 PID} & Tuer un processus \\
\cmd{df -h}      & Espace disque \\
\cmd{du -sh dossier} & Taille d'un dossier \\
\cmd{tar -czf archive.tar.gz dossier} & Archiver + compresser \\
\bottomrule
\end{tabularx}
\end{tcolorbox}

% ============================================================
\begin{tcolorbox}[basebox, colframe=orange, title=3. Gestion des permissions (chmod)]
Format : \cmd{-[u][g][o]} \quad ex: \cmd{-rwxr-xr--}\\[4pt]

\begin{tabular}{lll}
\toprule
\textbf{Valeur} & \textbf{Lettre} & \textbf{Droit} \\
\midrule
4 & r & Lecture (read) \\
2 & w & Écriture (write) \\
1 & x & Exécution (execute) \\
0 & - & Aucun \\
\bottomrule
\end{tabular}
\vspace{4pt}\\
\textbf{Exemples :}
\begin{itemize}
  \item \cmd{chmod 755 fichier} $\to$ \cmd{rwxr-xr-x}
  \item \cmd{chmod 644 fichier} $\to$ \cmd{rw-r--r--}
  \item \cmd{chmod 600 secret.key} $\to$ \cmd{rw-------}
  \item \cmd{chmod u+x script.sh} $\to$ ajoute exécution à l'user
\end{itemize}
\vspace{4pt}
\warn{Règle de sécurité} : ne jamais donner \textbf{w+x} simultanément à others (permet injection et exécution de code).\\[4pt]
\key{Propriétaire} : \cmd{chown user:group fichier}\\
\cmd{chown -R www-data:www-data /var/www/html}
\end{tcolorbox}

% ============================================================
\begin{tcolorbox}[basebox, colframe=teal, title=4. Gestion des utilisateurs et groupes]
\begin{tabular}{ll}
\toprule
\textbf{Commande} & \textbf{Action} \\
\midrule
\cmd{useradd mael} & Créer utilisateur (minimal) \\
\cmd{adduser mael} & Créer utilisateur (interactif) \\
\cmd{passwd mael} & Modifier mot de passe \\
\cmd{usermod -aG sudo mael} & Ajouter au groupe sudo \\
\cmd{groupadd devs} & Créer un groupe \\
\cmd{id mael} & Infos UID/GID de l'utilisateur \\
\cmd{who / whoami} & Utilisateurs connectés \\
\cmd{su - mael} & Changer d'utilisateur \\
\cmd{sudo commande} & Exécuter en root \\
\cmd{visudo} & Éditer les droits sudo \\
\bottomrule
\end{tabular}
\vspace{4pt}\\
\key{/etc/passwd} : liste des utilisateurs \\
\key{/etc/shadow} : mots de passe (hachés) \\
\key{/etc/group} : liste des groupes
\end{tcolorbox}

% ============================================================
\begin{tcolorbox}[basebox, colframe=navy, title=5. Gestion des services (systemd)]
\begin{tabular}{ll}
\toprule
\textbf{Commande} & \textbf{Action} \\
\midrule
\cmd{systemctl start apache2} & Démarrer un service \\
\cmd{systemctl stop apache2} & Arrêter \\
\cmd{systemctl restart apache2} & Redémarrer \\
\cmd{systemctl enable apache2} & Démarrage auto au boot \\
\cmd{systemctl disable apache2} & Désactiver au boot \\
\cmd{systemctl status apache2} & État du service \\
\cmd{journalctl -u apache2} & Logs du service \\
\bottomrule
\end{tabular}
\vspace{4pt}\\
\key{Services courants} : \cmd{apache2}, \cmd{nginx}, \cmd{ssh}, \cmd{mysql}, \cmd{cron}
\end{tcolorbox}

% ============================================================
\begin{tcolorbox}[basebox, colframe=green, title=6. Gestion des paquets (Debian/Ubuntu)]
\begin{tabular}{ll}
\toprule
\cmd{apt update} & Mettre à jour la liste \\
\cmd{apt upgrade} & Mettre à jour les paquets \\
\cmd{apt install paquet} & Installer \\
\cmd{apt remove paquet} & Désinstaller \\
\cmd{apt purge paquet} & Désinstaller + configs \\
\cmd{apt search nom} & Rechercher un paquet \\
\cmd{dpkg -l | grep nom} & Lister paquets installés \\
\bottomrule
\end{tabular}
\vspace{4pt}\\
\key{/etc/apt/sources.list} : dépôts configurés
\end{tcolorbox}

% ============================================================
\begin{tcolorbox}[basebox, colframe=crimson, title=7. Réseau sous Linux]
\begin{tabular}{ll}
\toprule
\cmd{ip addr show} & Afficher les interfaces \\
\cmd{ip route} & Table de routage \\
\cmd{ping 8.8.8.8} & Test connectivité \\
\cmd{ss -tuln} & Ports ouverts (successeur netstat) \\
\cmd{curl -I http://site.com} & Tester une URL \\
\cmd{ssh user@ip} & Connexion SSH \\
\cmd{scp fichier user@ip:/dest} & Copie sécurisée \\
\bottomrule
\end{tabular}
\vspace{4pt}\\
\key{/etc/network/interfaces} : config réseau (Debian) \\
\key{/etc/hosts} : résolution locale de noms \\
\key{/etc/resolv.conf} : serveurs DNS
\end{tcolorbox}

% ============================================================
\begin{tcolorbox}[basebox, colframe=orange, title=8. Fichiers importants du système]
\begin{tabular}{ll}
\toprule
\textbf{Fichier/Dossier} & \textbf{Rôle} \\
\midrule
\cmd{/etc/} & Configuration système \\
\cmd{/var/log/} & Journaux système \\
\cmd{/home/user/} & Répertoire personnel \\
\cmd{/tmp/} & Fichiers temporaires \\
\cmd{/bin/ et /usr/bin/} & Binaires/exécutables \\
\cmd{/boot/} & Noyau et GRUB \\
\cmd{/dev/} & Périphériques \\
\cmd{/proc/} & Processus (pseudo-FS) \\
\bottomrule
\end{tabular}
\end{tcolorbox}

% ============================================================
\begin{tcolorbox}[basebox, colframe=teal, title=9. PowerShell — Administration Windows]
\begin{tabular}{ll}
\toprule
\cmd{Get-ChildItem} & Lister (équiv. \cmd{ls}) \\
\cmd{Set-Location} & Changer répertoire (cd) \\
\cmd{New-Item} & Créer fichier/dossier \\
\cmd{Remove-Item -Recurse} & Supprimer \\
\cmd{Copy-Item} & Copier \\
\cmd{Get-Process} & Lister processus \\
\cmd{Get-Service} & Lister services \\
\cmd{Start-Service nom} & Démarrer service \\
\cmd{Get-EventLog -LogName System} & Journaux Windows \\
\cmd{Get-NetIPAddress} & Config réseau \\
\bottomrule
\end{tabular}
\vspace{4pt}\\
\key{Alias} : \cmd{ls}, \cmd{cd}, \cmd{cp} fonctionnent aussi en PS \\
\key{Pipeline} : \cmd{Get-Process | Where \{...\}}
\end{tcolorbox}

% ============================================================
\begin{tcolorbox}[basebox, colframe=navy, title=10. Sauvegarde \& Restauration]
\key{Stratégie 3-2-1} :
\begin{itemize}
  \item 3 copies des données
  \item Sur 2 supports différents
  \item Dont 1 hors site (cloud, distant)
\end{itemize}
\vspace{4pt}
\textbf{Types de sauvegarde :}
\begin{itemize}
  \item \key{Complète} : tout sauvegarder (lente, sûre)
  \item \key{Incrémentale} : depuis la dernière sauvegarde
  \item \key{Différentielle} : depuis la dernière complète
\end{itemize}
\vspace{4pt}
\key{Commandes tar :}\\
\cmd{tar -czf backup.tar.gz /data} \quad (créer)\\
\cmd{tar -xzf backup.tar.gz} \quad (extraire)\\[4pt]
\key{rsync} : synchronisation de fichiers :\\
\cmd{rsync -avz /src user@srv:/dest}
\end{tcolorbox}

\end{multicols}

\vspace{4pt}
\begin{tcolorbox}[basebox, colframe=crimson, title=Points clés à retenir pour l'examen]
\begin{multicols}{2}
\begin{itemize}
  \item \cmd{chmod 755} = \cmd{rwxr-xr-x} (user: tout, group+others: lecture+exec)
  \item \cmd{chmod 644} = \cmd{rw-r--r--} (fichier web standard)
  \item \warn{Ne jamais faire rm -rf /} (supprime tout le système)
  \item \key{sudo} = exécuter une commande en tant que root
  \item \cmd{systemctl enable} = démarrage automatique au boot
  \item \cmd{apt update} avant \cmd{apt install} (toujours)
  \item SSH port 22 — connexion sécurisée chiffrée
  \item \cmd{grep -r "texte" /dossier} = recherche récursive
  \item Propriétaire fichier : \cmd{chown user:group fichier}
  \item \cmd{/etc/passwd} : comptes, \cmd{/etc/shadow} : mots de passe hachés
\end{itemize}
\end{multicols}
\end{tcolorbox}

\end{document}
